% 清华大学本科毕业论文模板 - 具身智能无本体数据采集与模仿学习方法研究
\documentclass[UTF8,a4paper,12pt,fontset=fandol]{ctexart}

% 页面设置
\usepackage{geometry}
\geometry{top=3cm,bottom=3cm,left=3cm,right=3cm}

% 数学支持
\usepackage{amsmath,amssymb,amsfonts}

% 图表支持
\usepackage{graphicx}
\usepackage{booktabs}
\usepackage{multirow}
\usepackage{array}

% 列表环境
\usepackage{enumitem}

% 超链接
\usepackage{hyperref}
\hypersetup{colorlinks=true,linkcolor=black,citecolor=black}

% 参考文献
\usepackage[numbers,sort&compress]{natbib}

% 代码高亮（如果需要）
\usepackage{listings}
\lstset{basicstyle=\small\ttfamily}

% 算法伪代码
\usepackage{algorithm}
\usepackage{algorithmic}

% TikZ绘图
\usepackage{tikz}
\usetikzlibrary{shapes.geometric, arrows.meta, positioning, calc, backgrounds}

% 其他常用包
\usepackage{caption}
\usepackage{subcaption}
\usepackage{float}
\usepackage{xcolor}

% 标题格式设置
\ctexset{
    section = {
        format = \zihao{3}\heiti\centering,
        beforeskip = 24pt,
        afterskip = 18pt
    },
    subsection = {
        format = \zihao{4}\heiti,
        beforeskip = 24pt,
        afterskip = 6pt
    },
    subsubsection = {
        format = \zihao{-4}\heiti,
        beforeskip = 12pt,
        afterskip = 6pt
    }
}

% 行距
\linespread{1.5}

% 参考文献标题
\renewcommand{\bibname}{参考文献}

\begin{document}

% 封面
% 清华大学本科毕业论文封面
\thispagestyle{empty}
\begin{center}

\vspace*{2cm}

{\zihao{-0}\heiti 清华大学}

\vspace{0.5cm}

{\zihao{2}\heiti 本科生综合论文训练}

\vspace{2cm}

{\zihao{2}\heiti 具身智能无本体采集\\[0.3cm]模仿学习方法研究}

\vspace{3cm}

\begin{tabular}{rl}
\zihao{4}\heiti 院\hspace{2em}系：&\zihao{4} 电子工程系 \\[0.8cm]
\zihao{4}\heiti 专\hspace{2em}业：&\zihao{4} 电子信息科学与技术 \\[0.8cm]
\zihao{4}\heiti 学生姓名：&\zihao{4} 陈帅文 \\[0.8cm]
\zihao{4}\heiti 学\hspace{2em}号：&\zihao{4} 2022010701 \\[0.8cm]
\zihao{4}\heiti 指导教师：&\zihao{4} XXX教授 \\[0.8cm]
\end{tabular}

\vfill

{\zihao{4} 2026年6月}

\end{center}

\newpage

% 声明页
\thispagestyle{empty}
\begin{center}
{\zihao{3}\heiti 关于学位论文使用授权的说明}
\end{center}

\vspace{1.5cm}

\zihao{4}

本人完全了解清华大学有关保留、使用学位论文的规定，即：

清华大学拥有在著作权法规定范围内学位论文的使用权，其中包括：
（1）已获学位的研究生必须按学校规定提交学位论文，学校可以采用影印、
缩印或其他复制手段保存研究生上交的学位论文；（2）为教学和科研目的，
学校可以将公开的学位论文作为资料在图书馆、资料室等场所供校内师生阅读，
或在校园网上供校内师生浏览部分内容；（3）根据《中华人民共和国学位条例
暂行实施办法》，向国家图书馆报送可以公开的学位论文。

本人保证遵守上述规定。

\vspace{2cm}

\begin{flushright}
（保密的论文在解密后应遵守此规定）

\vspace{1cm}

论文作者签名：\underline{\hspace{4cm}}\hspace{2em}日期：\underline{\hspace{3cm}}

\vspace{0.5cm}

指导教师签名：\underline{\hspace{4cm}}\hspace{2em}日期：\underline{\hspace{3cm}}
\end{flushright}

\newpage

% 摘要页（中文）
\setcounter{page}{1}
\pagenumbering{Roman}

\begin{center}
{\zihao{3}\heiti 摘\hspace{1em}要}
\end{center}

\vspace{0.5cm}

\zihao{-4}

具身智能的发展高度依赖于大规模、高质量的机器人交互数据，但传统的遥操作采集方法依赖昂贵的机器人本体，难以满足数据规模化的需求。无本体数据采集方案提供了一种低成本、可扩展的替代路径，但其在数据采集、质量控制和策略训练等方面仍面临挑战。

本研究围绕"无本体的数据能否有效用于模仿学习"这一核心问题，构建了一套完整的无本体双臂采集-训练-部署系统，解决了三个关键技术挑战：

（1）视角矛盾问题：针对长程任务的全局视野需求，我们系统评估了多种第三视角相机方案，最终采用D435固定方案与DJI Nano头戴方案双线并进。针对Nano无线录制的同步难题，创新性地设计了双音chirp标定同步方案，实现$<$33ms的同步精度。

（2）噪声干扰问题：针对数据质量问题，建立了完整的预处理流水线（格式转换、降采样、时间对齐），设计了10条CHECK规则的数据筛选机制，并实施了轨迹相似性分析、HSIC依赖性分析和VLM表征分析等多种质量验证方法。

（3）信息嵌入问题：针对无本体数据缺乏固定坐标系的挑战，创新性地设计了Relative Action表示方案，解决无本体采集没有固定baselink的核心难题，使无本体数据能够有效用于策略训练。

实验验证了系统的有效性。在单臂pick-place任务上，400条数据达到90\%成功率，接近Franka遥操作的100\%，验证了跨平台泛化能力。反直觉地发现纯图像输入方案（90\%）优于图像+显式状态输入（75\%）。以handover任务作为双臂验证场景，初步验证了双臂协调的可行性，并分析了抓取失败、对齐失败、接取失败三类主要失败模式。

\vspace{0.5cm}
\noindent\textbf{关键词：}具身智能；无本体采集；模仿学习；双臂协调；数据质量控制

\newpage

% Abstract（英文）
\begin{center}
{\zihao{3}\heiti Abstract}
\end{center}

\vspace{0.5cm}

\zihao{-4}

The development of embodied AI heavily relies on large-scale, high-quality robot interaction data. However, traditional teleoperation collection methods depend on expensive robot hardware, making it difficult to meet the demands of data scaling. Embodiment-free data collection offers a low-cost, scalable alternative, but still faces challenges in data acquisition, quality control, and policy training.

This study addresses the core question of "Can embodiment-free data be effectively used for imitation learning?" by constructing a complete embodiment-free bimanual data collection-training-deployment system, solving three key technical challenges:

(1) Perspective conflict: To meet the global vision requirement of long-range tasks, we systematically evaluate multiple third-view camera solutions, adopting both D435 fixed and DJI Nano head-mounted options. For Nano's wireless recording synchronization challenge, we innovatively design a dual-tone chirp calibration synchronization scheme, achieving $<$33ms synchronization accuracy.

(2) Noise interference: Regarding data quality issues, we establish a complete preprocessing pipeline (format conversion, downsampling, time alignment), design a 10-CHECK-rule data filtering mechanism, and implement multiple quality verification methods including trajectory similarity analysis, HSIC dependence analysis, and VLM representation analysis.

(3) Information embedding: Addressing the challenge of embodiment-free data lacking fixed coordinate systems, we innovatively design the Relative Action representation scheme, solving the core problem of embodiment-free collection without fixed baselink, enabling effective use of embodiment-free data for policy training.

Experiments validate the system's effectiveness. On the single-arm pick-place task, 400 demonstrations achieve 90\% success rate, close to the 100\% of Franka teleoperation baseline, demonstrating cross-platform generalization. Counter-intuitively, we find that pure image input (90\%) outperforms image with explicit state input (75\%). Using the handover task as a bimanual validation scenario, we preliminary verify the feasibility of bimanual coordination and analyze three main failure modes: grasping failure, alignment failure, and receiving failure.

\vspace{0.5cm}
\noindent\textbf{Keywords:} Embodied AI; Embodiment-free Collection; Imitation Learning; Bimanual Coordination; Data Quality Control

\newpage


% 目录
\tableofcontents
\newpage

% 正文
% 第1章 绪论
\section{绪论}

\subsection{研究背景与意义}

近年来，具身智能（Embodied AI）受到学术界和工业界的广泛关注\cite{duan2022embodied}。与传统的"认知智能"不同，具身智能强调智能体与物理环境的实时交互能力，要求机器人能够通过感知、决策和执行的闭环，在真实世界中完成复杂任务\cite{billard2019learning,kroemer2021review}。这一领域已被纳入国家"十五五"规划建议，成为需要着力打造的新经济增长点，在工业制造、家庭服务、医疗康复、仓储物流等领域拥有巨大的潜在市场。

具身智能的进步同样遵循数据驱动的基本规律。观察认知智能领域的发展历程可以发现，模型性能的提升与训练数据规模呈正相关关系——当语言模型的训练数据从1T增长到十余T时，其理解与生成能力实现了质的飞跃。类似地，机器人学习系统的性能也高度依赖于交互数据的规模与多样性。公开文献显示，早期机器人模型使用数十万条轨迹训练，而近期研究通过引入互联网图文数据实现了更好的泛化能力，也有项目通过百万级异构数据训练获得了通用操作能力。这些数据密集型研究共同指向一个结论：具身智能系统的竞争力取决于数据获取的规模化和多样性保障。

然而，构建大规模的具身数据集面临着现实约束。传统的数据获取方式主要依赖人工遥操作机器人完成任务，这种方案存在多重局限：当采用主从机械臂架构时，需要配置与执行端同构的操控端设备，导致硬件投入成本高昂，且采集场景受限于机械臂的工作空间，此外操作者观察视角与机器人传感器视角的差异可能引入数据偏差；当采用虚拟现实设备遥操作时，虽然降低了对物理机械臂的依赖，但操作者缺乏真实的力触觉反馈，且系统存在网络传输延迟和位姿估计误差，这些因素都会降低操作精度。综合来看，基于机器人本体的传统采集方式在成本、场景适应性和扩展性方面难以满足规模化数据生产的需求。

为突破上述局限，学术界和工业界开始探索不依赖机器人本体的数据获取新范式。这类方法的核心思想是：数据采集者仅通过穿戴或手持专门的传感设备，在真实环境中完成操作演示并记录多模态数据，后续将这些演示数据用于训练控制策略。由于无需机器人本体参与采集过程，这种方法显著降低了数据生产的硬件门槛，使得规模化数据获取成为可能。近年来，多个研究团队推出了相应的采集系统，这些系统通常集成了手持操作端、可穿戴相机和惯性测量单元等组件，形成了轻量化的数据采集解决方案。

本研究旨在构建一套完整的无本体双臂遥操作采集系统，覆盖从数据采集、处理、训练到部署的全流程，并通过跨平台对比实验验证其有效性。研究的主要意义在于：为具身智能研究提供一种低成本、可扩展的数据获取路径，降低数据采集的硬件门槛，推动更大规模具身智能模型的训练与应用落地。

\subsection{研究内容与目标}

本研究围绕无本体数据采集与模仿学习展开，主要研究内容包括以下几个方面：

\textbf{（1）无本体采集系统构建}

设计并实现一套完整的无本体双臂数据采集系统，包括硬件平台选型、软件系统实现、数据格式转换与同步机制。硬件方面采用FastUMI/XV双臂采集设备配合固定位置D435第三视角相机；软件方面实现rosbag到mp4/json再到LeRobot格式的完整数据流程。

\textbf{（2）双臂协调策略研究}

针对双臂协调的挑战，研究有效的数据采集与策略学习方法。探索纯图像输入方案在双臂任务中的可行性，分析放弃显式状态表示（state input）的原因，并以handover任务作为双臂协调的关键验证场景。

\textbf{（3）数据质量分析方法论}

建立数据质量评估体系，设计10条CHECK规则用于数据筛选，实现轨迹相似性分析（基于DTW距离和层次聚类）和图像-动作依赖性分析（基于HSIC互信息估计），为数据采集提供质量监控工具。

\textbf{（4）跨平台实验验证}

通过多组对比实验验证系统有效性：单臂端到端验证（分析数据量-性能关系）、跨平台对比（FastUMI vs XV vs Franka遥操）、第三视角效果分析、双臂handover探索性实验。

本研究的目标是构建一套完整的无本体双臂采集-训练-部署系统，验证其在pick-place等单臂任务和handover等双臂任务上的可行性，为低成本具身数据采集提供可复用的技术方案。

\subsection{论文组织结构}

本文共分为七章，各章内容安排如下：

第一章为绪论，介绍研究背景、意义、内容与目标，以及论文的组织结构。

第二章为相关工作，综述无本体采集方法、双臂机器人遥操作技术、模仿学习策略以及数据质量评估方法的研究现状。

第三章为硬件系统与数据采集，解决"长程任务需求与采集设备头部视角缺失的矛盾"，详细介绍第三视角相机选型（D435固定方案与DJI Nano头戴方案）、多传感器数据同步（含创新的双音chirp标定方案），以及系统实现与性能指标。

第四章为数据处理与质量验证，解决"噪声数据的干扰问题"，阐述数据预处理关键技术（格式转换、降采样、时间对齐）、数据筛选机制，以及轨迹相似性和图像-动作依赖性分析的质量验证方法。

第五章为模型训练与信息嵌入，解决"无本体数据如何实现信息嵌入"，探讨输入设计（纯图像方案）、状态表示探索（ArUco码方案尝试）、核心创新点——Action表示选择（Relative Action vs Absolute Action），以及训练部署流程。

第六章为实验验证与Pipeline体现，展示完整pipeline在单臂任务（pick-place，数据量-性能关系、跨平台对比）和双臂任务（handover，失败模式分析）上的验证结果。

第七章为结论与展望，总结全文工作，指出主要贡献与创新点，并展望后续研究方向。

% 第2章 相关工作
\section{相关工作}

本章综述与本研究相关的技术背景和研究现状，包括无本体数据采集方法\cite{chi2024umi,wang2024dexcap}、双臂机器人遥操作技术\cite{zhao2023aloha,gello2023}、模仿学习策略\cite{argall2009survey,osa2018algorithmic}以及数据质量评估方法\cite{deminf2025,gretton2005hsic}四个方面的内容。

\subsection{无本体数据采集方法}

无本体数据采集是具身智能领域近年来的重要研究方向，其核心思想是在不依赖真实机器人硬件的情况下获取高质量的演示数据。

\textbf{（1）基于手持末端的数据采集}

通用操作接口（Universal Manipulation Interface）是这一类采集方案的典型代表。Chi等人\cite{chi2024umi}设计的系统让操作者手持专门的操作端完成演示，通过配置在腕部的广角相机捕捉视觉信息，利用侧向反光镜辅助估计深度，同时结合惯性测量单元记录手部运动。这种架构使得在真实场景中获取操作演示数据无需机器人本体在场，为模仿学习提供了数据来源。

FastUMI在UMI的基础上进行了硬件和算法优化，提高了数据采集设备的可靠性和易用性。本研究采用的FastUMI系统继承了这一设计思路，使用XR机器人进行双臂数据采集，进一步降低了硬件成本。

\textbf{（2）基于动作捕捉的数据采集}

Wang等人\cite{wang2024dexcap}开发的DexCap方案融合了动作捕捉技术，通过数据手套记录手指的精细运动，结合腕部视觉和惯性传感器追踪手臂轨迹。该方案在UMI基础上增设了胸前广角相机，增强了环境感知的稳定性，适用于需要灵巧操作的复杂任务场景。

\textbf{（3）众包数据采集}

Generalist AI团队开发了指套式采集设备，通过大众力量众包采集机器人演示数据。据报道，该方案峰值时每周可获取超过一万小时的人类演示数据。Sunday Robotics也采用类似方案专注于家庭服务机器人领域。这些工作证明无本体采集配合众包模式可以突破传统数据匮乏的瓶颈。

\subsection{双臂机器人遥操作技术}

双臂协调是机器人操作中的重要课题，遥操作技术为双臂数据采集提供了直接手段。

\textbf{（1）主从机械臂遥操作}

经典的主从遥控制架构要求操作者使用与执行端结构相同的主控设备，通过运动学映射实现动作复现。ALOHA系统\cite{zhao2023aloha}验证了这种架构在家庭操作任务中的可行性，GELLO框架\cite{gello2023}则探索了低成本硬件适配方案。这类方案能够采集高精度的运动轨迹，但设备投入较大，需要同时维护两套机械臂系统。

\textbf{（2）VR遥操作}

Open TeleVision系统\cite{cheng2024opentelevision}通过虚拟现实界面提供沉浸式的操作体验，利用手柄或手套追踪人体动作并实时驱动机器人运动。国内也有团队采用类似技术路线构建了大规模操作数据集。这类方案的优势在于无需操作者配备机械臂本体，但存在操作延迟和缺少触觉反馈的局限。

\textbf{（3）无本体双臂采集}

本研究的无本体双臂采集结合了上述两类方法的优点：既无需昂贵的机械臂本体（降低成本），又能获得人直接操作的精细动作（保证数据质量），同时避免了VR系统的延迟问题。

\subsection{模仿学习策略}

模仿学习（Imitation Learning）是机器人从人类演示中学习策略的主要方法。

\textbf{（1）行为克隆}

行为克隆（Behavior Cloning）是最基础的模仿学习方法\cite{argall2009survey,osa2018algorithmic}，通过监督学习直接从演示数据中学习状态到动作的映射。该方法简单有效，但在分布外数据上容易累积误差\cite{osa2018algorithmic}。

\textbf{（2）扩散策略}

Diffusion Policy\cite{chi2023diffusion}将机器人动作生成建模为去噪扩散过程，能够表达多模态分布并生成平滑轨迹。RDP（Reactive Diffusion Policy）\cite{xue2025rdp}采用"慢-快"双策略架构，以视觉主导产生规划，以触觉反馈快速调整接触细节，在需要精细力控的任务上表现优异。

\textbf{（3）视觉-语言-动作模型}

VLA（Vision-Language-Action）模型将预训练的视觉-语言模型接入机器人控制。Google的RT-2\cite{rt22023}利用海量图像-文本训练数据，实现指令驱动的操作和新技能零样本泛化。本研究采用的π0.5模型延续了这一架构，通过大规模预训练提升策略泛化能力。

\subsection{数据质量评估方法}

数据质量直接影响模仿学习的最终性能，近年来研究者提出了多种评估方法。

\textbf{（1）轨迹分析方法}

轨迹相似性分析使用动态时间规整（DTW）距离度量不同轨迹间的相似程度，结合层次聚类可识别数据中的主要模式。该方法能够发现离群演示，评估数据覆盖度。轨迹大数据综述文献讨论了这些方法在机器人数据挖掘中的应用。

\textbf{（2）互信息估计}

状态-动作互信息可用于衡量数据的可预测性和多样性。Robot Data Curation with Mutual Information Estimators\cite{deminf2025}利用互信息估计进行机器人演示数据的策展。希尔伯特-施密特独立性准则（HSIC）为互信息估计提供了计算框架。

\textbf{（3）分布距离度量}

最大均值差异（MMD）和推土机距离（Wasserstein Distance）可用于度量不同批次数据间的分布差异。MMD Imitation Learning\cite{mmdil2013}利用嵌入分布距离刻画演示与策略差异。本研究的轨迹相似性分析和图像-动作依赖性分析借鉴了这些方法。

综上所述，已有研究为本课题提供了重要基础：无本体采集设备为数据规模化提供了可能，VLA模型和扩散策略为策略学习提供了有效工具，数据质量分析方法为评估采集数据提供了手段。本文在这些基础上，针对双臂协调和系统完整性展开深入研究。

% 第3章 硬件系统与数据采集
\section{硬件系统与数据采集}

\subsection{引言：长程任务的视角需求}

在双臂机器人操作中，长程任务（如双手交接物品、跨桌面协作）需要操作者对环境全局有充分感知。然而，传统无本体采集方案仅依赖腕部相机，只能捕捉手部附近的局部视野，难以满足长程任务的全局视野需求。本章针对"长程任务需求与采集设备头部视角缺失的矛盾"，系统介绍了第三视角相机的选型探索、多传感器数据同步方案以及系统实现。

\subsection{第三视角相机选型}

\subsubsection{为什么需要第三视角}

腕部相机虽然能够捕捉精细的手部操作，但存在以下局限：

\begin{itemize}[leftmargin=2em]
    \item \textbf{视野局限}：仅能看到手部附近区域，无法感知双臂相对位置和全局环境
    \item \textbf{长程任务困难}：在handover等需要双手协调的任务中，腕部相机无法同时观察两只手
    \item \textbf{空间推理缺失}：策略难以建立全局空间地图，影响复杂操作规划
\end{itemize}

第三视角相机提供了全局视野，使策略能够观察双臂协作过程、物体位置和操作空间布局，是解决长程任务视角需求的关键。

\subsubsection{候选方案评估}

我们对四种第三视角方案进行了系统评估，评估维度包括FOV覆盖率、视野稳定性、舒适度和同步便利性。

\begin{table}[htbp]
\centering
\caption{第三视角相机选型对比}
\label{tab:camera_comparison}
\begin{tabular}{lccccc}
\toprule
方案 & FOV & 稳定性 & 舒适度 & 同步 & 综合评分 \\
\midrule
D435固定（胸口/固定） & 4/5 & 5/5 & N/A & 5/5 & \textbf{4.7/5} \\
DJI Nano头戴 & 5/5 & 4/5 & 5/5 & 2/5 & 4.0/5 \\
Lumos头戴 & 5/5 & 4/5 & 3/5 & 4/5 & 4.0/5 \\
ZED头戴 & 4/5 & 3/5 & 2/5 & 4/5 & 3.3/5 \\
\bottomrule
\end{tabular}
\end{table}

\textbf{D435固定方案}（主线方案）：
\begin{itemize}[leftmargin=2em]
    \item \textbf{优势}：ROS同步成熟可靠、视角稳定无抖动、部署便利、技术文档完善
    \item \textbf{部署方式}：固定在三脚架上，放置于操作台前方，高度约1.2-1.5米
    \item \textbf{局限}：佩戴位置固定，无法跟随头部视角；视角范围受固定位置限制
    \item \textbf{适用场景}：需要稳定全局视野的固定场景任务，如桌面操作、货架整理
\end{itemize}

\textbf{DJI Nano头戴方案}（创新探索）：
\begin{itemize}[leftmargin=2em]
    \item \textbf{优势}：
    \begin{itemize}
        \item 轻便舒适（仅50g，远轻于Lumos的200g+）
        \item FOV覆盖率5/5，广角镜头覆盖双臂活动范围
        \item 有配套头戴支架，佩戴几乎无感
        \item 无线录制，无需布线，操作者活动自由
    \end{itemize}
    \item \textbf{挑战}：同步问题（最初是主要瓶颈，需要外部事件触发）
    \item \textbf{解决}：开发了双音chirp标定同步方案（详见3.2.2节）
    \item \textbf{数据量}：40分钟录制约13GB，可接受
    \item \textbf{适用场景}：需要跟随视角的移动场景，或操作者需要频繁移动的任务
\end{itemize}

\textbf{Lumos头戴方案}（备选）：
\begin{itemize}[leftmargin=2em]
    \item \textbf{优势}：鱼眼180度FOV，视野拉满；稳定性较高
    \item \textbf{局限}：相机较重（约200g），长时间佩戴疲劳；数据传输冗余大，带宽要求高
    \item \textbf{综合评价}：视野优秀，但重量和带宽是硬伤
\end{itemize}

\textbf{ZED头戴方案}（放弃）：
\begin{itemize}[leftmargin=2em]
    \item \textbf{问题}：无配套头戴支架，需要DIY固定；相机本身较重，头部佩戴负担大
    \item \textbf{综合评价}：重量和体积导致佩戴体验差，不适合头戴场景
\end{itemize}

\subsubsection{方案选择与应用}

最终我们采用\textbf{双线并进}策略：
\begin{itemize}[leftmargin=2em]
    \item \textbf{D435固定}作为主线方案，用于大多数固定场景采集（如拿放杯子、上货任务）
    \item \textbf{DJI Nano头戴}作为创新方案，同步问题已解决，可用于特殊场景或需要跟随视角的任务
\end{itemize}

这种双方案设计为无本体采集提供了灵活性和可靠性保障。

\subsection{多传感器数据同步}

时间同步是确保数据可用性的关键环节。我们针对D435固定方案和Nano头戴方案分别设计了同步方案。

\subsubsection{D435 ROS时间同步}

D435相机\cite{intel2018realsense}原生支持ROS\cite{quigley2009ros,koubaa2016ros}，采用硬件时间戳实现精确同步。

\textbf{技术实现}：
\begin{enumerate}[leftmargin=2em]
    \item \textbf{硬件时间戳}：D435和FastUMI/XV设备在同一ROS网络中，使用统一系统时钟
    \item \textbf{时间对齐}：以较短时间序列为基准，对较长序列进行线性插值
    \item \textbf{误差控制}：确保左右臂每一帧时间戳差异小于33ms（对应30fps帧间隔）
    \item \textbf{异常标记}：超过阈值的片段标记为同步异常，参与后续筛选
\end{enumerate}

\textbf{同步精度}：实测平均误差约15ms，满足双臂协调任务要求。

\textbf{实现细节}：
\begin{itemize}[leftmargin=2em]
    \item ROS网络配置：所有设备连接同一WiFi或有线网络
    \item 时间同步服务：使用NTP或chrony确保系统时钟一致
    \item 数据录制：使用rosbag record同时录制所有topic
    \item 后期处理：rosbag解析时提取时间戳，进行对齐验证
\end{itemize}

\subsubsection{Nano头戴相机同步探索}

针对DJI Nano无线录制场景，传统ROS同步不适用，我们开发了\textbf{双音chirp标定同步方案}。

\textbf{问题背景}：
\begin{itemize}[leftmargin=2em]
    \item Nano无线录制，无法通过ROS获取时间戳
    \item 需要与主系统（FastUMI/XV）精确对齐
    \item 后期手动对齐费时费力且准确性低
    \item 早期尝试使用语音识别（ASR）进行对齐，但ASR在静音片段和重叠语音上表现不佳
\end{itemize}

\textbf{解决方案：双音chirp标定法}：

采用两种可区分的高频短音标记开始/结束\cite{oppenheim1999discrete,rabiner1975chirp}：
\begin{itemize}[leftmargin=2em]
    \item \textbf{开始音（Start）}：上调频chirp，2.2kHz $\rightarrow$ 4.4kHz，时长0.12秒
    \item \textbf{结束音（Stop）}：下调频chirp，4.5kHz $\rightarrow$ 2.2kHz，时长0.12秒
    \item \textbf{采样率}：48kHz（与视频录音一致）
    \item \textbf{频率选择}：选择2-5kHz范围，避开常见环境噪声频段，同时低于8kHz的语音截止频率
\end{itemize}

\textbf{检测算法}：
\begin{enumerate}[leftmargin=2em]
    \item \textbf{信号检测}：FFT归一化互相关（NCC）检测chirp信号
    \item \textbf{时间戳提取}：从音轨中提取精确的Start/Stop时间点
    \item \textbf{成对配对}：贪心算法将Start和Stop配对，期望起-停-起-停交替模式
    \item \textbf{误差容差}：$\pm$0.5秒内配对，支持未配对标记再扫描
\end{enumerate}

\textbf{技术细节}：
\begin{itemize}[leftmargin=2em]
    \item \textbf{FFT长度}：使用1024点FFT，频率分辨率约47Hz（48kHz采样率）
    \item \textbf{检测阈值}：NCC峰值超过0.8认为检测到有效信号
    \item \textbf{后处理流程}：抽音轨 $\rightarrow$ 模板检测（FFT互相关）$\rightarrow$ 成对配对 $\rightarrow$ ffmpeg切片
\end{itemize}

\textbf{同步精度}：实测误差$<$33ms，满足双臂协调要求。

\textbf{方案优势}：
\begin{itemize}[leftmargin=2em]
    \item \textbf{不依赖ASR}：完全绕过语音识别，避免语义模型弱点
    \item \textbf{物理波形定位}：时间戳由物理波形产生，比口语字级更稳定、可复现
    \item \textbf{实时轻量}：播标与检标共用同一份参考WAV，便于互相关匹配
    \item \textbf{可人工复核}：提供浏览器标注工具，支持人工调整和确认
\end{itemize}

这一方案创新性地解决了头戴相机同步难题，为无线录制场景提供了可靠的时间对齐方法。

\subsection{系统实现与性能}

\subsubsection{端到端数据流程}

系统涵盖从数据采集到策略部署的完整流程，共四个阶段：

\textbf{阶段一：数据采集}：
\begin{itemize}[leftmargin=2em]
    \item 启动ROS节点（XV SDK或Franka控制器）
    \item 启动FastUMI/XV采集监控工具
    \item 启动第三视角相机（D435或Nano）
    \item 使用rosbag record录制数据
    \item 每条轨迹录制6-10秒，对应30-60帧（60Hz采样，目标10Hz）
\end{itemize}

\textbf{阶段二：本地处理}：
\begin{itemize}[leftmargin=2em]
    \item 将rosbag转换为mp4+json格式（视频压缩+轨迹提取）
    \item 可视化检查每条轨迹，剔除明显异常数据
    \item 对数据进行排序和重新编号
    \item 压缩数据并上传到服务器
\end{itemize}

\textbf{阶段三：服务器处理}：
\begin{itemize}[leftmargin=2em]
    \item 解压数据并转换为LeRobot格式
    \item 计算归一化统计量（mean, std, q01, q99）
    \item 降采样至目标帧率（10Hz）
    \item 数据筛选（应用10条CHECK规则）
    \item 合并多批次数据并采样（100/200/400条等配置）
\end{itemize}

\textbf{阶段四：训练与部署}：
\begin{itemize}[leftmargin=2em]
    \item 基于π0.5模型微调（20,000步，约11-13小时）
    \item 模型权重下载到本地4090台式机
    \item 启动策略服务（WebSocket服务）
    \item 机械臂执行评测
\end{itemize}

\subsubsection{系统性能指标}

\begin{table}[htbp]
\centering
\caption{系统性能指标}
\label{tab:system_performance}
\begin{tabular}{lcc}
\toprule
指标 & 数值 & 说明 \\
\midrule
单条数据处理时间 & 约2-3分钟 & rosbag到LeRobot（本地） \\
服务器处理时间 & 约5-10分钟 & 解压+格式转换+归一化 \\
每小时数据存储量 & 约13GB（含视频） & 双臂+第三视角 \\
数据压缩比 & 约10:1 & rosbag到mp4 \\
同步精度（D435） & $<$33ms & ROS时间戳 \\
同步精度（Nano） & $<$33ms & chirp标定 \\
单模型训练时间 & 11-13小时 & 20k步，batch=16 \\
模型推理速度 & 10-15fps & 实时控制 \\
\bottomrule
\end{tabular}
\end{table}

\subsubsection{关键问题与解决}

\textbf{问题一：数据采集过热}
\begin{itemize}[leftmargin=2em]
    \item 现象：采集设备过热导致帧率下降，6秒视频只有70+帧（正常应有360帧）
    \item 解决：增加散热措施，控制单次采集时长，设备冷却后再次采集
\end{itemize}

\textbf{问题二：初始位置固定}
\begin{itemize}[leftmargin=2em]
    \item 现象：无本体采集每次初始位置可能不同，导致策略泛化困难
    \item 解决：使用脚踏板固定操作者位置，确保每次采集初始位置一致
\end{itemize}

\textbf{问题三：数据分布差异}
\begin{itemize}[leftmargin=2em]
    \item 现象：不同采集者操作速度和风格不同（位移步长差1.43倍，旋转步长差2.2倍）
    \item 解决：统一采集规范，增加数据多样性，后续通过数据筛选剔除异常分布数据
\end{itemize}

\subsubsection{对比平台}

为验证无本体方案的有效性，我们同时构建了\textbf{Franka遥操平台}作为对比基线。该平台采用Gello框架实现双臂遥操作，硬件配置包括：

\begin{itemize}[leftmargin=2em]
    \item \textbf{机械臂}：Franka Panda双臂（7DOF+7DOF）
    \item \textbf{夹爪}：Robotiq 2F-85（遥操）/ UMI夹爪（无本体）
    \item \textbf{相机}：腕部鱼眼相机（遥操）/ 手持设备相机（无本体）
    \item \textbf{第三视角}：D435固定位置相机（两者共用）
\end{itemize}

该平台用于第6章的跨平台对比实验，验证无本体数据与遥操数据的性能差异。


% 第4章 数据处理与质量验证
\section{数据处理与质量验证}

\subsection{引言：噪声数据的挑战}

原始采集的数据不可避免地包含异常片段，如操作失误、同步错误、设备故障等。这些"噪声数据"如果不加处理直接用于训练，会干扰策略学习，降低最终性能。本章针对"模仿学习中噪声数据的干扰问题"，系统介绍数据预处理技术、筛选机制和质量验证方法。

\subsection{数据预处理}

数据预处理是实现采集到训练闭环的基础环节，包含格式转换、降采样和时间对齐等关键技术。

\subsubsection{格式转换：rosbag到mp4+json}

原始数据以rosbag格式存储，包含ROS topic的时间序列数据。为降低存储开销和提升后续处理效率，我们将其转换为mp4视频+json轨迹的格式。

\textbf{转换流程}：
\begin{enumerate}[leftmargin=2em]
    \item \textbf{图像编码}：读取图像topic，编码为mp4视频文件（H.264编码，质量预设medium）
    \item \textbf{轨迹提取}：读取夹爪状态和位姿topic，保存为json格式
    \item \textbf{时间戳保留}：提取时间戳信息，建立图像帧与状态记录的对应关系
    \item \textbf{元数据记录}：记录采集者、任务类型、设备信息、原始帧率等
\end{enumerate}

\textbf{存储优化效果}：
\begin{itemize}[leftmargin=2em]
    \item \textbf{压缩比}：rosbag到mp4+json的压缩比约为10:1
    \item \textbf{存储节省}：每小时原始数据约130GB（rosbag，60Hz双目+RGB），转换后约13GB
    \item \textbf{处理效率}：mp4格式便于视频解码，提升训练时数据加载速度（使用imageio迭代器流式读取）
    \item \textbf{可视化友好}：mp4可直接播放，便于人工检查数据质量
\end{itemize}

\textbf{实现细节}：
\begin{itemize}[leftmargin=2em]
    \item 使用OpenCV或ffmpeg进行视频编码
    \item json格式包含：records数组（每帧的pose、clamp、timestamp）、fps、source等
    \item pose字段：7维[x, y, z, qx, qy, qz, qw]（位置+四元数）
    \item clamp字段：夹爪宽度（mm）或脉冲信号（遥操数据）
\end{itemize}

\subsubsection{降采样：提升学习效率}

原始采集数据通常为60fps（FastUMI/XV设备），但策略学习不需要如此高的时间分辨率。过高的帧率会引入冗余信息，增加计算负担，甚至导致策略难以收敛。

\textbf{降采样策略}：
\begin{itemize}[leftmargin=2em]
    \item \textbf{目标帧率}：从60fps降至10fps（每6帧取1帧）
    \item \textbf{降采样方法}：均匀采样，保持时间间隔一致
    \item \textbf{stride计算}：$stride = round(orig\_fps / target\_fps)$，实际帧率 $achieved\_fps = orig\_fps / stride$
\end{itemize}

\textbf{效果分析}：
\begin{itemize}[leftmargin=2em]
    \item 数据量减少83\%（从60fps到10fps），显著降低存储和计算成本
    \item 去除冗余帧，策略更容易学习关键动作（如抓取、放置的临界点）
    \item 保持足够的时间精度（100ms间隔），满足操作需求
    \item 训练时batch\_size可设置更大（16-128），提升训练效率
\end{itemize}

\textbf{实验验证}：
\begin{itemize}[leftmargin=2em]
    \item 对比10Hz和20Hz：10Hz在拿放杯子任务上效果更好，模型更容易收敛
    \item 过高帧率（20Hz+）会导致模型学习过多细微动作，反而降低泛化能力
\end{itemize}

\subsubsection{时间对齐}

双臂数据的时间对齐是确保协调性的关键步骤。

\textbf{对齐流程}：
\begin{enumerate}[leftmargin=2em]
    \item 以较短时间序列为基准（避免插值引入误差）
    \item 对较长序列进行线性插值或截断
    \item 确保左右臂每一帧时间戳差异小于33ms（对应30fps帧间隔）
    \item 标记同步误差过大的片段（$>$50ms）
\end{enumerate}

\textbf{异常处理}：
\begin{itemize}[leftmargin=2em]
    \item 时间戳差异$>$50ms的帧标记为同步异常
    \item 同步异常率$>$5\%的episode建议剔除
    \item 对于Nano数据，使用chirp标定的时间戳进行硬对齐
\end{itemize}

\subsection{数据筛选机制}

\subsubsection{筛选规则设计}

基于实验过程中的经验总结和大量调试，我们设计了10条CHECK规则用于数据质量评估。这些规则分为三个严重级别：

\begin{table}[htbp]
\centering
\caption{数据质量筛选规则体系}
\label{tab:check_rules}
\begin{tabular}{clp{5.5cm}cc}
\toprule
规则编号 & 检测维度 & 异常判定条件 & 严重程度 & 优先级 \\
\midrule
R01 & 文件完备性 & 缺少关键数据文件（视频/轨迹/元数据） & 致命 & 最高 \\
R02 & 序列长度 & 帧数低于正常范围的30\%（如$<$20帧） & 警告 & 中等 \\
R03 & 格式一致性 & 同批次数据分辨率/格式不统一 & 警告 & 中等 \\
R04 & 时间同步 & 单臂数据缺失或双臂时差$>$50ms & 致命 & 高 \\
R05 & 运动停滞 & 连续多帧无明显位移（$<$0.001m）或高比例静止 & 警告 & 中等 \\
R06 & 位移异常 & 单步移动距离$>$0.1m或旋转角度$>$0.5rad & 警告 & 中等 \\
R07 & 轨迹退化 & 整体运动范围过小（全程位移$<$0.05m） & 警告 & 中等 \\
R08 & 数据缺失 & 单臂有效帧数极少（$<$10帧） & 致命 & 最高 \\
R09 & 终点跳变 & 轨迹末尾出现非预期大位移（$>$0.05m） & 警告 & 中等 \\
R10 & 执行器状态 & 夹爪开合动作不充分（全程变化$<$10mm） & 轻微 & 低 \\
\bottomrule
\end{tabular}
\end{table}

\textbf{规则设计思路}：
\begin{itemize}[leftmargin=2em]
    \item \textbf{致命级别}：直接导致数据不可用的问题（文件缺失、同步失败）
    \item \textbf{警告级别}：可能影响训练效果的问题（运动异常、轨迹退化）
    \item \textbf{轻微级别}：对训练影响较小的问题（夹爪开合不充分）
\end{itemize}

\subsubsection{筛选流程与实施}

筛选流程自动化执行上述规则：

\begin{enumerate}[leftmargin=2em]
    \item \textbf{批量扫描}：遍历数据集所有episode，逐项执行CHECK规则
    \item \textbf{分级标记}：根据规则级别标记数据（致命/警告/轻微）
    \item \textbf{自动剔除}：致命级别数据自动剔除，不参与训练
    \item \textbf{人工复核}：警告级别数据生成HTML可视化报告供人工复核
\end{enumerate}

\textbf{人工复核工具}：
\begin{itemize}[leftmargin=2em]
    \item 浏览器端可视化界面，显示每条轨迹的视频和关键统计量
    \item 支持标记"删除"、"保留"、"待定"
    \item 批量操作支持（同批次相似异常可一键处理）
\end{itemize}

\textbf{工程实践价值}：

本研究建立了数据筛选流程并应用于实际数据采集。筛选机制能够识别明显的录制异常（如设备掉线、同步失败等），在handover数据集上剔除约6\%的离群数据。对于微妙的数据质量问题（如操作不够流畅），规则筛选的效果有限，需要结合第4.3节的可视化分析方法进一步评估。

\subsection{数据质量可视化分析}

为多角度验证数据质量，我们实施了轨迹相似性分析和图像-动作依赖性分析。这些方法不仅用于筛选，也用于理解数据分布和指导采集策略。

\subsubsection{轨迹相似性分析}

轨迹相似性分析用于识别数据中的主要操作模式和离群轨迹。

\textbf{方法}：
\begin{itemize}[leftmargin=2em]
    \item \textbf{DTW距离}\cite{berndt1994dtw,mueller2007dtw}：动态时间规整距离，衡量两条轨迹的相似程度，允许时间轴的非线性对齐
    \item \textbf{层次聚类}\cite{schubert2017clustering}：基于DTW距离矩阵，使用凝聚式聚类识别主要操作模式
    \item \textbf{离群检测}：识别与主流模式差异过大的轨迹（距离$>$3倍中位数）
\end{itemize}

\textbf{实施结果}（handover任务数据集，约80条演示轨迹）：
\begin{itemize}[leftmargin=2em]
    \item \textbf{模式识别}：识别出3个主要轨迹模式簇
    \begin{itemize}
        \item 簇A（约45\%）：标准的左递右接模式，轨迹平滑
        \item 簇B（约35\%）：双手同时移动的对称模式，同步性好
        \item 簇C（约20\%）：其他变体或需要更多步骤的复杂模式
    \end{itemize}
    \item \textbf{离群检测}：识别出5条异常轨迹（约6\%）
    \begin{itemize}
        \item 原因：操作失误（中途掉落物品）、录制问题（设备抖动）、非标准操作
        \item 处理：经人工复核后剔除，避免噪声干扰训练
    \end{itemize}
\end{itemize}

\textbf{实际指导价值}：
\begin{itemize}[leftmargin=2em]
    \item 识别出数据分布的主要模式，指导后续采集应侧重哪些风格
    \item 发现离群数据，及时剔除避免干扰训练
    \item 验证数据多样性，确保覆盖不同操作方式
\end{itemize}

\subsubsection{图像-动作依赖性分析}

图像-动作依赖性分析用于衡量视觉观测与执行动作之间的统计关联强度，评估数据的可学习性。

\textbf{方法}：
\begin{itemize}[leftmargin=2em]
    \item \textbf{HSIC}\cite{gretton2005hsic}：希尔伯特-施密特独立性准则，衡量两个变量（图像特征、动作特征）的统计依赖性
    \item \textbf{特征提取}：
    \begin{itemize}
        \item 图像特征：使用SigLIP\cite{siglip2023}预训练模型提取，768维
        \item 动作特征：PCA降维至10维（保留95\%方差）
    \end{itemize}
    \item \textbf{结果解读}：HSIC值越高，表示图像与动作关联越强，策略越容易学习
\end{itemize}

\textbf{实施结果}：

\begin{table}[htbp]
\centering
\caption{不同任务类型的图像-动作依赖性分析结果}
\label{tab:hsic_results}
\begin{tabular}{lccc}
\toprule
数据集（任务类型） & HSIC值 & 样本数 & 分析 \\
\midrule
bagging\_0430（放置） & 0.00456 & 512 & 依赖性较低 \\
handover\_umi\_0429（交接） & 0.02451 & 512 & 依赖性显著 \\
duck\_0430（抓取） & 0.01865 & 512 & 依赖性中等 \\
\bottomrule
\end{tabular}
\end{table}

\textbf{关键发现与指导意义}：
\begin{itemize}[leftmargin=2em]
    \item \textbf{任务类型差异显著}：交接类任务（handover，HSIC=0.025）的依赖性明显高于放置类任务（bagging，HSIC=0.005），说明交接任务中策略需要从视觉输入中提取更多空间关系信息
    \item \textbf{数据需求解释}：较高的HSIC值意味着策略需要学习更复杂的映射关系，这也解释了为何handover任务需要比单臂任务更多的数据支持（600条 vs 200条）
    \item \textbf{数据质量指标}：HSIC值可作为数据质量评估的参考，过低（$<$0.001）可能表示"僵死"数据（动作与图像无关）
    \item \textbf{采集策略指导}：对于HSIC值高的任务，应注重采集多样性；对于HSIC值低的任务，应注重动作一致性
\end{itemize}

\subsubsection{VLM表征空间分析}

为深入理解不同数据源（遥操 vs 无本体）和不同任务类型在模型内部表征空间中的分布特性，我们开展了基于VLM（Vision-Language Model）的表征分析。该方法参考Physical Intelligence《Emergence of Human to Robot Transfer》附录C的做法，通过分析微调后模型的内部表征，揭示数据域之间的相似性与差异性。

\textbf{分析方法}：
\begin{itemize}[leftmargin=2em]
    \item \textbf{特征提取}：使用微调后的π0.5模型，提取VLM（PaliGemma language\_model）最后一层输出
    \item \textbf{池化策略}：对prefix前K个token（K≈200）做mean-pool，得到2048维表征向量
    \item \textbf{降维可视化}：使用PCA、t-SNE、UMAP将高维表征降至2D进行可视化
    \item \textbf{分离度度量}：计算不同数据组质心间的L2距离，量化域间差异
\end{itemize}

\textbf{实验设置}：
我们对比了以下数据组在VLM表征空间中的分布（每组64个样本）：
\begin{itemize}[leftmargin=2em]
    \item \textbf{umi\_model\_robot\_data}：使用UMI微调模型前向提取的遥操数据表征
    \item \textbf{umi\_model\_umi\_data}：使用UMI微调模型前向提取的无本体handover数据表征
    \item \textbf{umi\_model\_umi\_other\_data}：使用UMI微调模型前向提取的无本体duck（抓取）数据表征
    \item \textbf{umi\_model\_umi\_bagging\_data}：使用UMI微调模型前向提取的无本体bagging（放置）数据表征
\end{itemize}

\textbf{表征分离度分析结果}：

\begin{table}[htbp]
\centering
\caption{VLM表征空间中的组间分离度（质心L2距离）}
\label{tab:vlm_separation}
\begin{tabular}{lc}
\toprule
数据组对比 & 质心L2距离 \\
\midrule
遥操数据 vs UMI handover数据 & 2.79 \\
遥操数据 vs UMI抓取数据 & 7.09 \\
遥操数据 vs UMI放置数据 & 5.20 \\
UMI handover vs UMI抓取 & 7.18 \\
UMI handover vs UMI放置 & 5.86 \\
UMI抓取 vs UMI放置 & 3.31 \\
\bottomrule
\end{tabular}
\end{table}

\textbf{组内分布分析}：
\begin{table}[htbp]
\centering
\caption{各数据组内部到质心的平均距离（表征紧凑性）}
\label{tab:vlm_compactness}
\begin{tabular}{lcc}
\toprule
数据组 & 平均距离 & 标准差 \\
\midrule
遥操数据 & 3.90 & 1.04 \\
UMI handover数据 & 6.29 & 2.06 \\
UMI抓取数据 & 4.24 & 1.70 \\
UMI放置数据 & 4.07 & 1.58 \\
\bottomrule
\end{tabular}
\end{table}

\textbf{关键发现}：
\begin{enumerate}[leftmargin=2em]
    \item \textbf{遥操与UMI数据存在明显分离}：即使使用相同的UMI微调模型，遥操数据和UMI无本体数据在表征空间中仍呈现不同的分布区域（质心距离2.79），说明两类数据采集域的固有差异在模型内部仍有体现
    
    \item \textbf{不同任务类型表征差异显著}：
    \begin{itemize}
        \item 抓取任务（duck）与handover任务分离度最高（7.18），说明两类任务在视觉特征上差异最大
        \item 放置任务（bagging）与handover任务分离度中等（5.86），体现了不同操作类型的内在差异
        \item 抓取与放置任务相对接近（3.31），均为单臂操作，视觉场景相似
    \end{itemize}
    
    \item \textbf{遥操数据表征更紧凑}：遥操数据组内到质心的平均距离（3.90）明显小于UMI handover数据（6.29），标准差也更小（1.04 vs 2.06），说明遥操数据采集更标准化、一致性更高
    
    \item \textbf{UMI数据分布更分散}：UMI handover数据的组内距离和方差都较大，反映了无本体采集过程中操作者动作风格的多样性，这也解释了为何需要更多数据量来达到稳定性能
\end{enumerate}

\textbf{PCA方差解释}：
\begin{itemize}[leftmargin=2em]
    \item 第一主成分解释64.2\%的方差
    \item 第二主成分解释7.2\%的方差
    \item 前两维累计解释71.4\%的方差，说明2D可视化能够较好地反映原始高维空间的结构
\end{itemize}

\textbf{可视化结果解读}：
\begin{itemize}[leftmargin=2em]
    \item t-SNE和UMAP降维结果均显示，不同数据组形成了可区分的聚类
    \item 遥操数据点分布相对集中，呈紧凑团状
    \item UMI无本体数据分布较分散，不同任务（抓取、放置、交接）形成不同区域
    \item 这种分离模式为理解模型的域适应性提供了直观依据
\end{itemize}

\textbf{指导意义}：
\begin{itemize}[leftmargin=2em]
    \item \textbf{数据混合策略}：表征空间的可分离性支持了混合遥操和UMI数据进行训练的策略，模型能够学习到跨域的共享表征
    \item \textbf{任务分类洞察}：不同任务类型在表征空间中的分布差异，为后续的多任务学习提供了特征层面的依据
    \item \textbf{数据质量评估}：表征紧凑性可作为数据一致性的量化指标，辅助筛选高质量数据子集
\end{itemize}

\subsubsection{多视角（三视角）嵌入一致性分析}

除了上述基于轨迹和图像-动作依赖性的分析方法外，我们还开发了多视角嵌入一致性分析方法，用于评估同一时刻三路相机（左腕、右腕、第三视角）在语义嵌入空间中的一致性，以及单路相机在时间序列上的连续性。

\textbf{分析动机}：
\begin{itemize}[leftmargin=2em]
    \item 三路相机（左腕、右腕、第三视角）应该观察到同一场景的不同视角
    \item 如果某路相机出现遮挡、黑屏、时间戳错位或极端运动模糊，其嵌入表征会出现异常
    \item 通过度量三路嵌入的相互一致性和单路时序连续性，可以识别潜在的采集质量问题
\end{itemize}

\textbf{技术方法}：
\begin{enumerate}[leftmargin=2em]
    \item \textbf{特征提取}：使用冻结的SigLIP编码器（与训练一致的图像预处理），分别提取左腕（$z_L$）、右腕（$z_R$）、第三视角（$z_B$）的嵌入向量
    \item \textbf{L2归一化}：对嵌入向量进行归一化，得到$\hat{z}_L, \hat{z}_R, \hat{z}_B$
    \item \textbf{成对余弦相似度}：计算三路嵌入之间的余弦相似度
    \begin{itemize}
        \item 左-右相似度：$c_{LR} = \hat{z}_L \cdot \hat{z}_R$
        \item 左-Base相似度：$c_{LB} = \hat{z}_L \cdot \hat{z}_B$
        \item 右-Base相似度：$c_{RB} = \hat{z}_R \cdot \hat{z}_B$
    \end{itemize}
    \item \textbf{单路时序连续性}（可选）：计算同一路相机相邻帧的余弦相似度
    \begin{itemize}
        \item $s_L = \hat{z}_{L,t} \cdot \hat{z}_{L,t+1}$（同一episode内）
        \item 同理计算$s_R, s_B$
    \end{itemize}
\end{enumerate}

\textbf{重要说明}：
\begin{itemize}[leftmargin=2em]
    \item 第三视角与手腕相机视场和内容差异大，\textbf{不得}期待两两余弦"普遍很高"
    \item 左腕-右腕（$c_{LR}$）通常内容相似（都拍摄手部操作），余弦值相对较高
    \item 左/右腕-第三视角（$c_{LB}, c_{RB}$）由于视场差异大，余弦绝对值常偏低，属预期
    \item 本分析依赖三种信号综合判断：
    \begin{enumerate}
        \item \textbf{批内相对分位}：与全局分布比较，识别异常低值
        \item \textbf{同episode内时间突刺}：某帧突然偏离该episode的median
        \item \textbf{单路时序连续性}：某路相机相邻帧突变，可能为遮挡或跳帧
    \end{enumerate}
\end{itemize}

\textbf{异常检测规则}：
\begin{enumerate}[leftmargin=2em]
    \item \textbf{Episode级异常}：若某episode的$\min(c_{LB})$低于全局p10分位数，或低于该episode内$median - k \cdot MAD$，则标记为异常
    \item \textbf{Frame级异常}：若某帧的$c_{LB}$或$s_B$的稳健Z-score（robust z-score）超过阈值（默认4.0），则标记为可疑帧
    \item \textbf{禁止}：使用单一绝对阈值（如"cos < 0.5即坏"）宣布好坏，因为不同相机对之间的基线相似度本就不同
\end{enumerate}

\textbf{输出指标}：
\begin{itemize}[leftmargin=2em]
    \item \textbf{逐帧指标}：$c_{LR}, c_{LB}, c_{RB}$（及可选的$s_L, s_R, s_B$）写入CSV
    \item \textbf{分位数统计}：各指标的全局p10/p50/p90
    \item \textbf{离群标记}：episode和frame级别的异常列表（JSON格式）
    \item \textbf{直方图}：三路成对余弦的分布可视化
\end{itemize}

\textbf{应用场景}：
\begin{enumerate}[leftmargin=2em]
    \item \textbf{标定验证}：检测三路相机时间戳是否对齐（若$c_{LR}$突然下降，可能为同步错位）
    \item \textbf{遮挡检测}：某路相机被遮挡时，与其他路的余弦相似度会异常下降
    \item \textbf{坏帧识别}：单路时序连续性$s$突变，可能为黑屏、运动模糊或丢帧
    \item \textbf{数据筛选补充}：作为10条CHECK规则的补充，从语义嵌入角度识别异常
\end{enumerate}

\textbf{与几何方法对比}：
\begin{itemize}[leftmargin=2em]
    \item \textbf{优势}：无需相机外参标定，纯语义嵌入计算，实现简单
    \item \textbf{局限}：弱于几何重投影误差（若已知相机标定），仅作为弱信号探针
    \item \textbf{定位}：不替代标定流程，而是作为数据质检的辅助工具
\end{itemize}

\subsubsection{数据分布可视化}

除上述分析方法外，我们还实现了多种数据分布可视化工具：

\textbf{Action分布直方图}：
\begin{itemize}[leftmargin=2em]
    \item 绘制各维度的分布直方图（位置x/y/z、旋转、夹爪）
    \item 对比不同批次数据的分布差异
    \item 发现分布偏移，指导归一化策略选择
\end{itemize}

\textbf{State分布可视化}：
\begin{itemize}[leftmargin=2em]
    \item 绘制TCP位姿的3D分布散点图
    \item 验证采集是否覆盖足够的工作空间
    \item 发现采集盲区，指导后续采集补全
\end{itemize}

\textbf{轨迹3D可视化}：
\begin{itemize}[leftmargin=2em]
    \item 绘制末端执行器在3D空间中的运动轨迹
    \item 叠加多条轨迹，观察一致性
    \item 标记异常轨迹，辅助人工复核
\end{itemize}

\subsection{归一化策略与数据增强}

\subsubsection{归一化方法对比}

在训练过程中，我们发现归一化方法对模型性能有显著影响。我们对比了两种归一化方法：

\textbf{Quantile归一化}：
\begin{equation}
x_{norm} = \frac{x - q_{01}}{q_{99} - q_{01} + \epsilon} \times 2 - 1
\end{equation}
将数据的1\%分位映射到-1，99\%分位映射到+1（线性缩放）。

\textbf{优点}：对heavy-tail/outlier更鲁棒，避免异常值影响整体分布。

\textbf{缺点}：
\begin{itemize}[leftmargin=2em]
    \item 改变数据分布形状，可能丢失重要信息
    \item 显式padding（用0填充）后经过quantile归一化会被拉到-1，导致模型学习偏差
    \item 在我们的实验中，使用quantile归一化导致部分数据维度（dim00、01、04、06）存在固定bias
\end{itemize}

\textbf{Z-score归一化}：
\begin{equation}
x_{norm} = \frac{x - \mu}{\sigma + \epsilon}
\end{equation}

\textbf{优点}：
\begin{itemize}[leftmargin=2em]
    \item 保留数据分布的相对关系，更适合模仿学习任务
    \item 显式padding的0在归一化后仍为0（或接近0），不会引入偏差
    \item 在我们的实验中，改用z-score归一化后，原本训练失败的数据能够正常训练
\end{itemize}

\textbf{最终选择}：采用z-score归一化作为默认归一化方法。

\subsubsection{数据增强策略}

虽然无本体采集的数据量已经较大（数百条），但我们仍探索了数据增强策略：

\textbf{颜色抖动}：
\begin{itemize}[leftmargin=2em]
    \item 对图像进行亮度、对比度、饱和度调整
    \item 模拟不同光照条件下的采集场景
\end{itemize}

\textbf{空间扰动}：
\begin{itemize}[leftmargin=2em]
    \item 对轨迹进行微小的时间平移（$\pm$2帧）
    \item 对初始位置进行随机扰动（但保持与图像一致）
\end{itemize}

\textbf{实际效果}：
\begin{itemize}[leftmargin=2em]
    \item 颜色抖动对拿放杯子任务帮助不大（背景固定）
    \item 初始位置扰动对泛化性有一定帮助，但需谨慎避免与图像不一致
    \item 总体而言，由于无本体数据量已经较大（200-400条），数据增强的收益有限
\end{itemize}

\subsection{本章小结}

本章建立了完整的数据处理与质量验证体系，包括：

\begin{enumerate}[leftmargin=2em]
    \item \textbf{预处理流程}：rosbag到mp4+json格式转换、10fps降采样、时间对齐
    \item \textbf{筛选机制}：10条CHECK规则自动检测异常，人工复核确保质量
    \item \textbf{质量分析}：
    \begin{itemize}
        \item 轨迹相似性分析识别模式簇和离群数据
        \item HSIC分析评估图像-动作依赖性
        \item VLM表征空间分析揭示遥操与UMI数据的域间差异（质心距离2.79）和不同任务类型的表征分离
        \item 多视角嵌入一致性分析从语义嵌入角度检测三路相机的同步异常、遮挡和坏帧
    \end{itemize}
    \item \textbf{归一化策略}：对比quantile和z-score归一化，最终采用z-score
\end{enumerate}

这些方法和工具不仅用于提升数据质量，也为理解数据分布、指导采集策略提供了量化依据。特别是：
\begin{itemize}[leftmargin=2em]
    \item VLM表征分析揭示了遥操数据表征更紧凑（组内距离3.90），UMI数据表征更分散（组内距离6.29），不同任务类型在表征空间中形成可区分的聚类（抓取vs handover分离度7.18），为跨域学习和多任务策略训练提供了特征层面的依据
    \item 多视角嵌入一致性分析通过度量左腕-右腕-第三视角三路嵌入的余弦相似度和时序连续性，为标定验证、遮挡检测和坏帧识别提供了语义层面的弱信号探针
\end{itemize}


% 第5章 模型训练与信息嵌入
\section{模型训练与信息嵌入}

\subsection{引言：无本体数据的信息嵌入挑战}

无本体采集的数据需要被模型有效理解和利用，才能训练出高性能的策略。本章针对"无本体数据如何实现信息嵌入"这一关键问题，探讨输入设计、状态表示和Action表示方案。

\subsection{输入设计}

\subsubsection{输入配置对比实验}

我们系统对比了两种输入配置的性能差异：纯图像输入 vs 图像+显式状态输入。这是本研究的重要发现之一。

\textbf{拿放杯子任务（400条无本体数据）}：

\begin{table}[htbp]
\centering
\caption{拿放杯子任务：不同输入模态的性能比较（400条演示数据）}
\label{tab:input_comparison}
\begin{tabular}{lcp{6cm}}
\toprule
策略输入配置 & 测试成功率 & 行为特征观察 \\
\midrule
视觉图像+绝对State & 75\% & 动作保守，末端轨迹抖动，偶尔出现明显偏差 \\
仅视觉图像 & \textbf{90\%} & 动作果断，末端轨迹平滑，泛化性好 \\
\bottomrule
\end{tabular}
\end{table}

\subsubsection{反直觉的发现与分析}

实验结果显示，纯图像输入方案（90\%）反而优于加入显式状态表示的方案（75\%）。这一发现初看反直觉，但可能有以下解释：

\textbf{模型自身推理能力}：
\begin{itemize}[leftmargin=2em]
    \item π0.5模型基于大规模视觉-语言预训练，具备强大的视觉推理能力
    \item 模型能够从图像中隐式提取空间位置信息、物体关系、操作进度
    \item 无需显式的State输入，模型自身可以"看懂"场景并做出决策
\end{itemize}

\textbf{State噪声问题}：
\begin{itemize}[leftmargin=2em]
    \item 无本体采集的State估计（如SLAM位姿）存在较大误差
    \item FastUMI/XV设备的SLAM在快速运动时容易丢帧
    \item State估计误差反而成为噪声，干扰策略学习
    \item 而图像信息相对完整，噪声更少
\end{itemize}

\textbf{端到端优势}：
\begin{itemize}[leftmargin=2em]
    \item 纯图像方案实现了真正的端到端学习，避免了中间表示的误差累积
    \item 策略直接从像素学习到动作，无需经过State估计的"翻译"
    \item 减少了系统复杂度，提升了可靠性
\end{itemize}

\textbf{遥操数据的对比}：

值得注意的是，在遥操数据中，State输入确实有提升（上货任务80\%→92.5\%）。这说明：
\begin{itemize}[leftmargin=2em]
    \item State输入的价值取决于State质量
    \item 遥操数据的State来自精确的机械臂编码器，质量高
    \item 无本体数据的State来自SLAM估计，质量较低
    \item 在State质量不足时，不如依赖视觉
\end{itemize}

\textbf{结论}：基于上述发现，本研究采用纯图像输入方案，简化了系统设计并获得了更好的性能。

\subsection{状态表示探索}

\subsubsection{ArUco码方案尝试}

在第三视角坐标系统一的探索中，我们曾考虑基于ArUco码的检测方案。该方案通过在场景中放置标定板，利用计算机视觉识别ArUco码来建立统一的世界坐标系。

\textbf{方案设计}：
\begin{itemize}[leftmargin=2em]
    \item 在操作空间中放置ArUco码标定板
    \item 使用第三视角相机检测ArUco码，计算相机坐标系到世界坐标系的变换
    \item 通过第一帧的双臂位姿，计算左右臂相对于世界坐标系的位置
    \item 后续帧使用相对位移，建立双臂统一的坐标表示
\end{itemize}

\textbf{方案局限}：
\begin{itemize}[leftmargin=2em]
    \item 需要额外的标定设备部署，增加了系统复杂度
    \item 动态场景中标定板可能被遮挡或移动
    \item 增加了部署成本和时间
    \item 对于单臂任务，坐标系统一并非必需
\end{itemize}

\textbf{最终决定}：综合考虑后，该方案未进入实际大规模测试阶段。我们最终采用纯图像输入方案，依靠模型自身的空间推理能力处理视角差异，避免了复杂的坐标系统一过程。

\subsection{Action表示选择（核心）}

\subsubsection{无本体采集的特殊挑战：缺少固定baselink}

\textbf{问题背景}：

在传统遥操作中，机器人有固定的基座（baselink），所有动作都相对于这个固定坐标系表示。然而，无本体采集中：
\begin{itemize}[leftmargin=2em]
    \item 操作者佩戴设备移动，没有固定的世界坐标系原点
    \item 每次采集的初始位置可能不同
    \item 绝对坐标表示会导致策略难以泛化
    \item 左右臂分别有自己的坐标系，缺乏统一参考
\end{itemize}

\subsubsection{Action表示方案对比}

我们系统对比了三种Action表示方案：

\textbf{Absolute Action（绝对动作）}：
\begin{itemize}[leftmargin=2em]
    \item 定义：相对于固定世界坐标系的绝对位姿
    \item 问题：无本体采集缺乏固定原点，绝对坐标在不同采集间差异大
    \item 结果：策略难以泛化，部署时效果很差（机械臂基本不动或动作异常）
\end{itemize}

\textbf{Delta Action（增量动作）}：
\begin{itemize}[leftmargin=2em]
    \item 定义：相对于上一时刻的位姿变化
    \item 优点：不依赖固定坐标系
    \item 缺点：累积误差，长序列容易漂移
    \item 结果：训练时loss下降，但部署效果仍不理想
\end{itemize}

\textbf{Relative Action（相对动作）- 最终方案}：
\begin{itemize}[leftmargin=2em]
    \item 定义：相对于当前episode第一帧的位姿变化（$\Delta$pose from start）
    \item 优点：
    \begin{itemize}
        \item 不依赖固定坐标系，适应无本体场景
        \item 表示的是"如何移动"而非"移动到哪里"，更具通用性
        \item 与纯图像输入方案配合，实现端到端学习
        \item 消除了baselink不一致的问题
    \end{itemize}
    \item 结果：训练稳定，泛化能力强，部署效果好
\end{itemize}

\begin{table}[htbp]
\centering
\caption{不同Action表示方案的性能对比（拿放杯子任务，100条遥操数据）}
\label{tab:action_representation}
\begin{tabular}{lcc}
\toprule
Action表示 & 成功率 & 备注 \\
\midrule
Absolute Action & $<$10\% & 几乎无法工作 \\
Delta Action & 30-50\% & 有一定效果但不稳定 \\
Relative Action & \textbf{100\%} & 稳定可靠 \\
\bottomrule
\end{tabular}
\end{table}

\subsubsection{具体实现}

参考UMI原文方法\cite{chi2024umi}和旋转表示相关研究\cite{saxena2008mathematical,zhou2019continuity}，我们采用以下Action表示：

\begin{equation}
Action_t = [\Delta x, \Delta y, \Delta z, rot_{6d}, gripper]
\end{equation}

其中：
\begin{itemize}[leftmargin=2em]
    \item $[\Delta x, \Delta y, \Delta z]$：末端位置相对于起始帧的三维位移
    \item $rot_{6d}$：末端姿态相对于起始帧的旋转变化，使用6D表示（两个3维向量）
    \item $gripper$：夹爪开合状态（归一化后的宽度）
\end{itemize}

\textbf{关键设计选择}：
\begin{enumerate}[leftmargin=2em]
    \item \textbf{相对位移}：使用相对于起始帧的增量，而非绝对坐标或上一帧增量
    \item \textbf{6D旋转表示}：采用6维旋转表示（参考Zhou等人2019\cite{zhou2019continuity}），避免万向锁问题，连续性更好
    \item \textbf{归一化}：对位移量进行z-score归一化，使其分布在合理范围（约$\pm$2）
    \item \textbf{Gripper处理}：遥操数据使用脉冲信号，无本体数据使用连续宽度值，统一归一化到$[-1, 1]$
\end{enumerate}

\textbf{实现细节}：
\begin{itemize}[leftmargin=2em]
    \item 每个episode的第一帧Action为0（相对于自身的位移为0）
    \item 使用LeRobot的action存储下一帧的状态，解决chunk第一帧为0的问题
    \item 归一化时跳过gripper维度（第10维），使用原始值
\end{itemize}

\subsubsection{方案有效性验证}

该Action表示方案在实验中表现出良好的效果：
\begin{itemize}[leftmargin=2em]
    \item \textbf{单臂任务}：拿放杯子任务400条数据达到90\%成功率
    \item \textbf{跨平台泛化}：FastUMI和XV平台均能使用相同表示
    \item \textbf{训练稳定性}：相比绝对坐标表示，训练过程更稳定，收敛更快
    \item \textbf{部署可靠性}：部署时机械臂执行流畅，无明显抖动或偏差
\end{itemize}

这一方案创新性地解决了无本体采集缺乏固定坐标系的问题，是无本体数据能够有效用于训练的关键技术之一。

\subsection{训练策略与部署}

\subsubsection{微调配置}

基于LeRobot数据集\cite{lerobot2024}，使用π0.5模型\cite{pi0_5_2025}进行微调：

\textbf{模型配置}：
\begin{itemize}[leftmargin=2em]
    \item \textbf{预训练模型}：π0.5\cite{pi0_2024}（大规模预训练的VLA模型）
    \item \textbf{输入}：纯图像（224×224分辨率，RGB）
    \item \textbf{输出}：Relative Action（10维：3维位移+6维旋转+1维夹爪）
    \item \textbf{Action维度}：32（pi0.5原生使用32维动作）
    \item \textbf{Action Horizon}：10（预测未来10步的动作）
\end{itemize}

\textbf{训练配置}：
\begin{itemize}[leftmargin=2em]
    \item \textbf{优化器}：AdamW
    \item \textbf{学习率}：5e-5（带warmup和cosine decay）
    \item \textbf{Warmup步数}：1,000步
    \item \textbf{训练步数}：20,000步
    \item \textbf{批次大小}：16-128（根据GPU资源和数据量调整）
    \item \textbf{保存间隔}：2,000步
    \item \textbf{训练时间}：约11-13小时（单卡A100/4090）
\end{itemize}

\textbf{归一化配置}：
\begin{itemize}[leftmargin=2em]
    \item 使用z-score归一化
    \item 归一化统计量（mean, std）从训练数据计算
    \item 跳过gripper维度的归一化（使用原始值）
\end{itemize}

\subsubsection{策略服务部署}

训练好的模型部署为实时策略服务：

\textbf{服务架构}：
\begin{itemize}[leftmargin=2em]
    \item 基于WebSocket的实时推理服务
    \item 输入：当前相机图像（RGB，224×224）
    \item 输出：预测的动作序列（10步Relative Action）
    \item 推理速度：约10-15fps，满足实时控制需求
\end{itemize}

\textbf{部署流程}：
\begin{enumerate}[leftmargin=2em]
    \item 从服务器下载训练好的模型权重（checkpoint）
    \item 在本地4090台式机启动策略服务（serve\_policy.py）
    \item 启动RealRL机器人控制服务（robot server）
    \item 启动评测客户端（eval\_policy.sh）
    \item 策略接收图像，输出动作，机械臂执行
\end{enumerate}

\textbf{关键参数}：
\begin{itemize}[leftmargin=2em]
    \item \textbf{控制频率}：10Hz（与训练时降采样频率一致）
    \item \textbf{动作平滑}：使用chunk内的动作序列，减少抖动
    \item \textbf{安全阈值}：设置动作幅值限制，防止机械臂碰撞
\end{itemize}

\subsubsection{训练过程中的问题与解决}

\textbf{问题一：Loss下降但部署效果差}
\begin{itemize}[leftmargin=2em]
    \item 现象：训练loss正常下降到0.08，但部署时机械臂动作异常
    \item 原因：归一化方法选择不当（quantile vs z-score）
    \item 解决：改用z-score归一化，效果明显改善
\end{itemize}

\textbf{问题二：Rotation表示选择}
\begin{itemize}[leftmargin=2em]
    \item 现象：使用3D轴角表示时，旋转维度学习效果差
    \item 原因：轴角表示存在不连续性问题（万向锁）
    \item 解决：改用6D旋转表示，训练更稳定
\end{itemize}

\textbf{问题三：Gripper控制}
\begin{itemize}[leftmargin=2em]
    \item 现象：夹爪开合时机不准确，提前松手或未及时闭合
    \item 原因：gripper维度的归一化方式与其他维度不同
    \item 解决：归一化时skip gripper维度，使用原始宽度值
\end{itemize}

\textbf{问题四：数据分布差异}
\begin{itemize}[leftmargin=2em]
    \item 现象：不同采集者数据训出模型效果差异大
    \item 原因：操作速度和风格不同导致数据分布偏移
    \item 解决：统一采集规范，增加数据筛选，混合多批次数据训练
\end{itemize}

\subsection{从采集到部署的闭环}

完整的流程形成闭环：

\begin{enumerate}[leftmargin=2em]
    \item \textbf{数据采集}：FastUMI/XV + 第三视角相机（D435或Nano）
    \item \textbf{本地处理}：rosbag→mp4+json，可视化检查，剔除异常
    \item \textbf{上传服务器}：压缩数据，scp上传到eva14服务器
    \item \textbf{服务器处理}：解压→LeRobot格式→计算归一化统计量→数据筛选
    \item \textbf{模型微调}：基于LeRobot数据，π0.5微调20,000步
    \item \textbf{模型下载}：checkpoint下载到本地4090
    \item \textbf{策略部署}：启动WebSocket策略服务
    \item \textbf{评测执行}：机械臂执行任务，记录成功率
\end{enumerate}

这一闭环流程经过多次迭代优化，现已稳定可靠，支持快速迭代实验。

\subsection{本章小结}

本章针对无本体数据的信息嵌入问题，主要贡献包括：

\begin{enumerate}[leftmargin=2em]
    \item \textbf{输入设计}：发现纯图像输入优于显式State输入的反直觉结论，简化了系统设计
    \item \textbf{状态表示}：探索了ArUco码方案，最终采用纯图像方案
    \item \textbf{Action表示}：创新性地提出Relative Action方案，解决无本体采集缺乏固定baselink的核心问题
    \item \textbf{训练部署}：建立了完整的训练-部署流程，解决了归一化、旋转表示、gripper控制等关键问题
\end{enumerate}

这些技术方案共同支撑了无本体数据的有效利用，是实现完整pipeline的关键。


% 第6章 实验验证与Pipeline体现
\section{实验验证与Pipeline体现}

本章展示完整pipeline在单臂和双臂任务中的验证结果，证明无本体数据采集与模仿学习方案的可行性。实验涵盖三类典型任务：拿放杯子（pick-place-cup）、上货（stock-shelves）、开笔记本（open-laptop），并对比了不同数据规模、不同输入配置下的性能表现。

\subsection{实验任务设计}

\subsubsection{任务一：拿放杯子（pick-place-cup）}

\textbf{任务描述}：机械臂需要从桌面指定区域抓取杯子，并放置到目标盘子中。该任务要求精确的抓取位姿控制和放置位置精度，是评估策略基本操作能力的标准任务。

\textbf{任务难度}：中等。涉及精细的抓取操作（杯子边缘较薄）和准确的放置定位（需要放入盘子而非桌面）。

\subsubsection{任务二：上货（stock-shelves）}

\textbf{任务描述}：机械臂需要将物品（如彩色方块）从桌面拿起，放置到货架的指定格子中。与pick-place任务相比，上货任务要求更高的空间定位和垂直方向的精确控制。

\textbf{任务难度}：中高。需要控制机械臂到达特定高度，并在三维空间中精确定位。评估发现，无本体数据训练的模型容易出现"提前放下"的问题，即尚未到达目标高度就松开夹爪。

\subsubsection{任务三：开笔记本（open-laptop）}

\textbf{任务描述}：机械臂需要打开合上的笔记本电脑。该任务包含多个子步骤：定位笔记本边缘、插入夹爪、掀开屏幕、按压确认。

\textbf{任务难度}：高。属于多阶段长程任务，要求策略具备较强的时序推理能力。实验表明，该任务对数据量要求最高，400条以下数据难以成功。

\subsection{单臂任务验证}

\subsubsection{实验设置}

\textbf{数据采集}：
\begin{itemize}[leftmargin=2em]
    \item \textbf{无本体数据}：使用FastUMI手持采集设备进行数据采集
    \item \textbf{遥操基线}：使用Gello遥操框架控制Franka机械臂采集数据
    \item \textbf{第三视角}：固定位置D435相机提供全局视野
    \item \textbf{数据格式}：rosbag原始数据，处理后转为LeRobot格式
\end{itemize}

\textbf{模型训练}：
\begin{itemize}[leftmargin=2em]
    \item \textbf{预训练模型}：π0.5（大规模VLA模型）
    \item \textbf{微调步数}：20,000步（约11-13小时训练时间）
    \item \textbf{批次大小}：16-128（根据GPU资源调整）
    \item \textbf{输入配置}：纯图像输入 / 图像+state输入
    \item \textbf{输出配置}：Relative Action（6D旋转+位置+gripper）
\end{itemize}

\textbf{评测方法}：
\begin{itemize}[leftmargin=2em]
    \item \textbf{测试次数}：每个模型测试20次
    \item \textbf{评价指标}：成功率（成功完成任务的次数占比）
    \item \textbf{位置泛化}：在多个位置（左、右、上、下）进行测试
    \item \textbf{ retry机制}：允许失败后重试，记录总尝试次数
\end{itemize}

\subsubsection{数据规模对性能的影响}

我们对拿放杯子任务进行了详细的数据规模-性能关系研究，对比了100条、200条、400条数据的训练效果。

\textbf{无本体数据（纯图像输入）}：

\begin{table}[htbp]
\centering
\caption{拿放杯子任务：无本体数据不同数据规模的性能对比（纯图像输入）}
\label{tab:data_scale_cup_embodiment_free}
\begin{tabular}{cccc}
\toprule
数据规模 & 成功率 & 左侧成功率 & 下方成功率 \\
\midrule
100条 & 77.5\% & 较好 & 较差 \\
200条 & 90\% & 90\% & 90\% \\
400条 & \textbf{90\%} & 90\% & 90\% \\
\bottomrule
\end{tabular}
\end{table}

\textbf{关键发现}：
\begin{itemize}[leftmargin=2em]
    \item \textbf{100条数据}：成功率77.5\%，能够完成基本任务但泛化性不足，对某些位置（如下方）表现较差
    \item \textbf{200条数据}：成功率跃升至90\%，达到实用水平，各位置泛化性能均衡
    \item \textbf{400条数据}：维持90\%成功率，与200条相比提升不明显，但模型更稳定
    \item \textbf{数据效率拐点}：200条左右是性价比最高的数据规模
\end{itemize}

\textbf{遥操基线数据（纯图像输入）}：

\begin{table}[htbp]
\centering
\caption{拿放杯子任务：遥操数据不同数据规模的性能对比（纯图像输入）}
\label{tab:data_scale_cup_teleop}
\begin{tabular}{cccc}
\toprule
数据规模 & 成功率 & 左侧成功率 & 下方成功率 \\
\midrule
50条 & 100\% & 100\% & 100\% \\
100条 & 100\% & 100\% & 100\% \\
\bottomrule
\end{tabular}
\end{table}

\textbf{关键发现}：
\begin{itemize}[leftmargin=2em]
    \item 遥操数据仅需50条即可达到100\%成功率
    \item 遥操数据的数据效率约为无本体数据的4倍
    \item 验证了遥操作为"完美基线"的有效性
\end{itemize}

\subsubsection{输入配置对比：纯图像 vs 图像+State}

我们对比了两种输入配置的性能差异，这是本研究的重要发现之一。

\textbf{拿放杯子任务（400条无本体数据）}：

\begin{table}[htbp]
\centering
\caption{拿放杯子任务：不同输入配置的性能对比}
\label{tab:input_config_cup}
\begin{tabular}{lcc}
\toprule
输入配置 & 成功率 & 备注 \\
\midrule
纯图像输入 & \textbf{90\%} & 动作果断，轨迹平滑 \\
图像+绝对State & 75\% & 动作保守，末端抖动 \\
\bottomrule
\end{tabular}
\end{table}

\textbf{上货任务（400条无本体数据）}：

\begin{table}[htbp]
\centering
\caption{上货任务：不同输入配置的性能对比}
\label{tab:input_config_shelves}
\begin{tabular}{lcc}
\toprule
输入配置 & 成功率 & 备注 \\
\midrule
纯图像输入 & 70\% & 有提前放下问题 \\
图像+绝对State & - & 未测试 \\
\bottomrule
\end{tabular}
\end{table}

\textbf{遥操数据上货任务（100条）}：

\begin{table}[htbp]
\centering
\caption{上货任务：遥操数据不同输入配置的性能对比}
\label{tab:input_config_shelves_teleop}
\begin{tabular}{lcc}
\toprule
输入配置 & 成功率 & 备注 \\
\midrule
纯图像输入 & 80\% & 有提前放下问题 \\
图像+绝对State & 92.5\% & 效果更好 \\
\bottomrule
\end{tabular}
\end{table}

\textbf{反直觉发现}：
\begin{itemize}[leftmargin=2em]
    \item 在拿放杯子任务中，\textbf{纯图像输入（90\%）反而优于图像+State（75\%）}
    \item 分析认为：
    \begin{itemize}
        \item π0.5模型具备强大的视觉推理能力，能够从图像中隐式提取空间信息
        \item 显式State估计可能引入噪声，反而干扰策略学习
        \item 无本体采集的State估计误差较大，不如依赖视觉
    \end{itemize}
    \item 在上货任务中，遥操数据的State输入确实有提升（80\%→92.5\%），说明State在复杂任务中有价值，但前提是State质量要足够高
\end{itemize}

\subsubsection{跨平台对比：遥操 vs 无本体}

我们系统对比了遥操（Franka+Gello）和无本体（FastUMI/XV）两种数据采集方式的性能。

\begin{table}[htbp]
\centering
\caption{不同采集平台的数据效率对比}
\label{tab:platform_comparison_detail}
\begin{tabular}{lcccc}
\toprule
任务 & 遥操50条 & 遥操100条 & 无本体100条 & 无本体400条 \\
\midrule
拿放杯子 & 100\% & 100\% & 77.5\% & 90\% \\
上货 & 67.5\% & 80\% & 62.5\% & 70\% \\
开笔记本 & 90\% & 100\% & 0\% & 50\% \\
\bottomrule
\end{tabular}
\end{table}

\textbf{数据效率分析}：
\begin{itemize}[leftmargin=2em]
    \item \textbf{拿放杯子}：遥操数据效率是无本体的约4倍（遥操50条=无本体200条）
    \item \textbf{上货}：遥操数据效率是无本体的约2-4倍
    \item \textbf{开笔记本}：无本体数据需要400条才达到50\%，而遥操100条即可达到100\%
    \item \textbf{结论}：无本体数据虽然效率较低，但通过增加数据量（4倍左右）可以达到接近遥操的性能
\end{itemize}

\subsubsection{开笔记本任务：高难度任务验证}

开笔记本是三个任务中难度最高的，需要多阶段推理和精细操作。

\textbf{实验结果}：

\begin{table}[htbp]
\centering
\caption{开笔记本任务：不同数据规模的性能}
\label{tab:open_laptop}
\begin{tabular}{lcc}
\toprule
数据规模 & 无本体成功率 & 遥操成功率 \\
\midrule
50条 & 0\% & 90\% \\
100条 & 0\% & 100\% \\
200条 & 0\% & - \\
300条 & 0\% & - \\
400条 & 50\% & - \\
\bottomrule
\end{tabular}
\end{table}

\textbf{关键发现}：
\begin{itemize}[leftmargin=2em]
    \item 该任务对数据量要求极高，无本体数据需要400条才能达到50\%成功率
    \item 遥操数据仅需50-100条即可达到90-100\%成功率
    \item 验证了复杂任务中无本体数据的可行性，但需要大规模数据采集
\end{itemize}

\subsection{双臂任务探索}

\subsubsection{handover任务设计}

\textbf{任务描述}：双臂机器人需要完成物品交接任务。具体流程为：
\begin{enumerate}[leftmargin=2em]
    \item 左臂从桌面抓取物品（如矿泉水瓶）
    \item 左臂将物品递送至双臂中间位置
    \item 右臂从中间位置接取物品
    \item 右臂将物品放置到右侧目标区域
\end{enumerate}

\textbf{任务难度}：高。涉及双臂协调、精确的空间对齐、时机配合。

\subsubsection{数据采集与训练}

\textbf{数据采集历程}：
\begin{itemize}[leftmargin=2em]
    \item \textbf{第一阶段}（3月）：使用D435第三视角相机，采集35条数据初步验证pipeline可行性
    \item \textbf{第二阶段}（4月）：完善数据质量，采用脚踏板固定初始位置，采集50余条数据
    \item \textbf{第三阶段}（4月中下旬）：大规模采集，总计约600条数据（包含不同批次混合）
    \item \textbf{第四阶段}（4月底）：使用DJI Nano作为头戴第三视角相机，采集约120条数据
\end{itemize}

\textbf{数据筛选}：
\begin{itemize}[leftmargin=2em]
    \item 使用10条CHECK规则筛选数据
    \item 剔除约6\%的明显异常数据（设备掉线、同步失败、操作失误等）
    \item 人工复核轨迹，保留高质量数据
\end{itemize}

\subsubsection{实验结果}

\textbf{训练配置对比}：

我们测试了多种训练配置，包括不同第三视角相机（D435固定 vs Nano头戴）、有无mask（是否遮盖第三视角）、不同数据规模等。

\begin{table}[htbp]
\centering
\caption{handover任务：不同配置的训练结果}
\label{tab:handover_results}
\begin{tabular}{llcc}
\toprule
配置编号 & 配置说明 & 数据量 & 效果评估 \\
\midrule
0322\_s1 & 高第三视角+Lumos相机 & 84条 & 效果很差，未学会 \\
0322\_s2 & 高第三视角+RealSense相机 & 84条 & 效果一般 \\
0322\_s3 & mask第三视角+Lumos & 84条 & 效果一般 \\
0322\_s4 & mask第三视角+RealSense & 84条 & 效果一般 \\
handover\_0329 & 重新采集高质量数据 & 50余条 & 效果一般 \\
handover\_0416\_mix & 混合多批次数据+筛选 & 约250条 & 效果较好 \\
handover\_0429 & 最终高质量数据 & 约600条 & 效果最好 \\
\bottomrule
\end{tabular}
\end{table}

\textbf{关键发现}：
\begin{itemize}[leftmargin=2em]
    \item \textbf{数据质量至关重要}：早期数据（0322批次）效果很差，经过筛选和重新采集后（0429批次）效果明显提升
    \item \textbf{数据规模}：600条混合数据效果最好，验证了大规模数据对双臂任务的重要性
    \item \textbf{第三视角价值}：高第三视角（相机位置较高）比mask（遮挡）效果更好，说明全局视野对双臂协调很重要
    \item \textbf{Nano相机可行性}：使用DJI Nano头戴相机采集的120条数据可以用于训练，验证了无线录制方案的可行性
\end{itemize}

\subsubsection{失败模式分析}

对handover任务失败案例进行详细分析，主要失败模式包括：

\begin{enumerate}[leftmargin=2em]
    \item \textbf{抓取失败（约35\%）}：
    \begin{itemize}
        \item 夹爪未完全闭合导致物品掉落
        \item 抓取位置偏移，物品在移动过程中滑落
        \item 对某些初始位置（随机接近位置）表现差，出现OOD（分布外）情况
    \end{itemize}
    
    \item \textbf{对齐失败（约30\%）}：
    \begin{itemize}
        \item 双手对齐位置不够精确，导致交接困难
        \item 高度不匹配，物品传递不顺畅
        \item 第三视角位置变化导致策略对空间位置估计偏差
    \end{itemize}
    
    \item \textbf{接取失败（约35\%）}：
    \begin{itemize}
        \item 右臂未能及时闭合夹爪接住物品
        \item 动作时机不协调，交接失败
        \item 递交过程中提前松手导致物品掉落
    \end{itemize}
\end{enumerate}

\subsection{训练过程中的关键问题与解决}

\subsubsection{归一化方法选择}

在训练过程中，我们发现归一化方法对模型性能有显著影响。

\textbf{问题发现}：
\begin{itemize}[leftmargin=2em]
    \item 使用quantile归一化（将1\%分位映射到-1，99\%分位映射到+1）时，部分数据训练效果很差
    \item 离线测试发现，某些数据维度（dim00、01、04、06）存在固定bias，dim08（gripper维度）未学到有效模式
\end{itemize}

\textbf{解决方案}：
\begin{itemize}[leftmargin=2em]
    \item 改用z-score归一化（zero-centered normalization）
    \item z-score公式：$x_{norm} = (x - \mu) / (\sigma + \epsilon)$
    \item 效果：改用z-score归一化后，原本训练失败的数据能够正常训练，offline测试结果明显改善
\end{itemize}

\textbf{原因分析}：
\begin{itemize}[leftmargin=2em]
    \item quantile归一化对重尾分布（heavy-tail）和异常值更鲁棒，但会改变数据分布形状
    \item z-score归一化保留数据分布的相对关系，更适合模仿学习任务
    \item 显式padding（用0填充）后经过quantile归一化会被拉到-1，导致模型学习偏差
\end{itemize}

\subsubsection{旋转表示选择}

我们对比了不同旋转表示方式的效果。

\begin{table}[htbp]
\centering
\caption{不同旋转表示方式的对比}
\label{tab:rotation_representation}
\begin{tabular}{llc}
\toprule
旋转表示 & 说明 & 效果 \\
\midrule
3D旋转（轴角） & 使用3维轴角表示 & 较差，容易学偏 \\
6D旋转 & 使用6维表示（两个3维向量） & \textbf{最好}，连续性更好 \\
\bottomrule
\end{tabular}
\end{table}

\textbf{结论}：采用6D旋转表示（参考Zhou等人2019的研究），避免万向锁问题，训练更稳定。

\subsubsection{数据分布问题}

我们发现不同采集者采集的数据分布存在差异，影响模型泛化能力。

\textbf{对比分析}（陈帅文200条 vs 龙哥100条）：
\begin{itemize}[leftmargin=2em]
    \item \textbf{轨迹长度}：258帧 vs 299帧
    \item \textbf{时间戳频率}：60.26Hz vs 60.47Hz
    \item \textbf{位移步长}（p99 median）：0.0131m vs 0.0092m（陈帅文快1.43倍）
    \item \textbf{旋转步长}（p99 median）：0.0269rad vs 0.0121rad（陈帅文快2.2倍）
\end{itemize}

\textbf{结论}：不同采集者的操作速度和风格存在差异，建议在采集时统一操作规范，或增加数据多样性以提升模型鲁棒性。

\subsection{实验总结}

\subsubsection{主要结论}

\begin{enumerate}[leftmargin=2em]
    \item \textbf{无本体数据可行}：通过增加数据量（约4倍），无本体数据可以达到接近遥操数据的性能
    \item \textbf{纯图像方案有效}：在简单任务中，纯图像输入反而优于显式State输入，说明π0.5模型具备强大的视觉推理能力
    \item \textbf{数据规模拐点}：拿放杯子任务200条数据达到90\%成功率，是上货和开笔记本任务的基础数据量
    \item \textbf{双臂任务可行但挑战性高}：handover任务需要600条以上高质量数据才能取得较好效果，且失败模式复杂
\end{enumerate}

\subsubsection{数据效率总结}

\begin{table}[htbp]
\centering
\caption{各任务数据效率总结}
\label{tab:data_efficiency_summary}
\begin{tabular}{lccc}
\toprule
任务 & 遥操达标数据量 & 无本体达标数据量 & 效率比 \\
\midrule
拿放杯子 & 50条（100\%） & 200条（90\%） & 4:1 \\
上货 & 100条（80\%） & 400条（70\%） & 4:1 \\
开笔记本 & 100条（100\%） & 400条（50\%） & 8:1 \\
handover（双臂） & 待验证 & 600条（较好） & - \\
\bottomrule
\end{tabular}
\end{table}

\subsubsection{未来改进方向}

\begin{itemize}[leftmargin=2em]
    \item \textbf{数据质量提升}：统一采集规范，增加数据筛选规则，提升单条数据质量
    \item \textbf{双臂协调优化}：探索更精细的双手协调表示，增加handover数据量至1000条以上
    \item \textbf{泛化能力}：测试跨场景、跨物体的泛化性能
    \item \textbf{数据效率}：研究如何通过数据增强、预训练等方法提升无本体数据效率
\end{itemize}


% 第7章 结论与展望
\section{结论与展望}

\subsection{工作总结}

本研究针对具身智能数据采集中的三个关键问题，构建了一套完整的无本体双臂采集-训练-部署系统，并通过实验验证了其可行性。

\subsubsection{解决视角矛盾：第三视角采集系统}

针对"长程任务需求与采集设备头部视角缺失的矛盾"，我们：
\begin{itemize}[leftmargin=2em]
    \item 系统评估了多种第三视角相机方案，最终采用\textbf{D435固定方案}作为主线，\textbf{DJI Nano头戴方案}作为创新探索
    \item 针对Nano无线录制的同步难题，创新性地设计了\textbf{双音chirp标定同步方案}，实现$<$33ms的同步精度
    \item 验证了第三视角相机在长程任务中的全局视野价值
\end{itemize}

\subsubsection{解决噪声干扰：数据处理与质量验证}

针对"模仿学习中噪声数据的干扰问题"，我们：
\begin{itemize}[leftmargin=2em]
    \item 建立了\textbf{预处理流水线}：rosbag→mp4+json格式转换（压缩比10:1）、10fps降采样、时间对齐
    \item 设计了\textbf{10条CHECK规则}的数据筛选机制，剔除约6\%的明显异常数据
    \item 实施了\textbf{多维度质量分析}：轨迹相似性分析识别3个主要模式簇和离群数据，HSIC分析揭示不同任务的依赖性差异
\end{itemize}

\subsubsection{解决信息嵌入：无本体数据的有效利用}

针对"无本体数据如何实现信息嵌入"，我们：
\begin{itemize}[leftmargin=2em]
    \item 验证了\textbf{纯图像输入方案}的有效性，发现纯图像（90\%）优于显式状态表示（75\%）
    \item 探索了状态表示方案（ArUco码尝试），最终采用简化的image-only方案
    \item 创新性地设计了\textbf{Relative Action表示}，解决无本体采集缺乏固定baselink的问题，是无本体数据能够有效用于训练的关键技术
\end{itemize}

\subsection{主要贡献}

本研究的主要贡献包括：

\begin{enumerate}[leftmargin=2em]
    \item \textbf{系统贡献}：构建了一套完整的无本体双臂采集-训练-部署系统，覆盖从数据采集、处理、训练到部署的全流程，降低了具身数据采集的硬件门槛
    
    \item \textbf{方法创新}：针对头戴相机同步难题，提出双音chirp标定同步方案，实现了无线录制场景下的精确时间对齐
    
    \item \textbf{技术突破}：设计Relative Action表示方案，创新性地解决无本体采集缺乏固定坐标系的问题，使无本体数据能够有效用于策略训练
    
    \item \textbf{实验验证}：通过多组对比实验验证了系统有效性，单臂pick-place任务达到90\%成功率，双臂handover任务初步验证可行
\end{enumerate}

\subsection{局限与展望}

\subsubsection{当前局限}

\begin{itemize}[leftmargin=2em]
    \item \textbf{数据量需求}：handover等双臂协调任务需要更多数据支持（与HSIC分析一致），当前80条数据难以达到理想的稳定性
    \item \textbf{任务复杂度}：系统在复杂精细操作上的性能仍有提升空间，handover任务成功率低于单臂任务
    \item \textbf{泛化能力}：策略的跨场景泛化能力需要进一步验证
\end{itemize}

\subsubsection{未来研究方向}

\begin{enumerate}[leftmargin=2em]
    \item \textbf{数据规模化}：扩大handover等复杂任务的数据采集规模，探索数据量-性能关系的边界
    \item \textbf{策略优化}：研究更精细的双手协调表示方法，提升双臂任务的稳定性和成功率
    \item \textbf{跨场景验证}：在更多样化的场景和任务中验证系统的泛化能力
    \item \textbf{人机协作}：探索人机协作采集模式，结合人工监督和自动化采集提升效率
    \item \textbf{实时性优化}：优化策略推理速度，实现更高频率的实时控制
\end{enumerate}

\subsection{结语}

本研究构建的无本体双臂采集-训练-部署系统，验证了无本体数据在具身智能研究中的可用性和潜力。通过解决视角矛盾、噪声干扰和信息嵌入三个关键问题，我们证明了一套低成本、可扩展的数据获取路径的可行性，为具身智能的规模化发展提供了技术基础。未来，随着数据采集规模的扩大和策略方法的优化，无本体采集有望在更多复杂场景中得到应用，推动具身智能技术的进一步发展。



% 参考文献
\bibliographystyle{gbt7714-numerical}
\bibliography{ref/refs}

% 致谢
% 致谢
\begin{center}
{\zihao{3}\heiti 致\quad 谢}
\end{center}
\addcontentsline{toc}{section}{致谢}

（致谢内容待补充）

\newpage


\end{document}
